% T3 — Текстовые задачи: концентрация / смеси / сплавы. ЕГЭ профиль №10. Класс 11.
\documentclass[a4paper,11pt]{article}
\usepackage{../../../../Lessons/_templates/preamble_worksheet}
\usepackage{../_shared/worksheet_check}

\renewcommand{\TaskBox}[2]{%
  \noindent\begin{tikzpicture}
    \node[draw=wcprimary, line width=0.8pt, rounded corners=4pt,
          inner xsep=3mm, inner ysep=2mm,
          text width=\dimexpr\linewidth-6mm\relax,
          anchor=north west] (B) at (0,0) {#2};
    \node[anchor=south west, fill=wcprimary, text=white,
          rounded corners=2pt, inner xsep=5pt, inner ysep=1.5pt,
          xshift=1.5pt] at (B.north west)
      {\scriptsize\textbf{Задача №#1}};
  \end{tikzpicture}\par\vspace{1.5mm}}

\begin{document}

\WorksheetTitle{Текстовые задачи. Концентрация и сплавы}
\WorksheetSubtitle{ЕГЭ профиль, задание №10 \quad·\quad Класс 11 \quad·\quad T3 \quad·\quad ID: T3-mix-2026-05-11}
\FirstPageHeader

\TaskBox{1}{%
  Имеется $100$~г сплава, содержащего $30\%$ серебра. Сколько граммов чистого серебра нужно добавить, чтобы получить сплав, содержащий $50\%$ серебра? Ответ дайте в граммах.\par
  \vspace{1.5mm}\hfill\textbf{\color{wcprimary}Ответ:}~\AnswerField{1}%
}

\TaskBox{2}{%
  Смешали $4$~кг $25\%$-го раствора и $6$~кг $35\%$-го раствора некоторого вещества. Сколько процентов составляет концентрация полученного раствора? Ответ дайте в процентах.\par
  \vspace{1.5mm}\hfill\textbf{\color{wcprimary}Ответ:}~\AnswerField{2}%
}

\TaskBox{3}{%
  Имеется два сплава меди и цинка. Первый кусок массой $80$~г содержит $40\%$ меди, второй массой $120$~г содержит $60\%$ меди. Сколько процентов меди будет в сплаве, полученном при их сплавлении? Ответ дайте в процентах.\par
  \vspace{1.5mm}\hfill\textbf{\color{wcprimary}Ответ:}~\AnswerField{3}%
}

\TaskBox{4}{%
  Из $8$~кг раствора, содержащего $25\%$ соли, выпарили $3$~кг воды. Сколько процентов соли содержится в полученном растворе? Ответ дайте в процентах.\par
  \vspace{1.5mm}\hfill\textbf{\color{wcprimary}Ответ:}~\AnswerField{4}%
}

\TaskBox{5}{%
  Виноград содержит $50\%$ воды, а изюм, получаемый из винограда, --- $25\%$ воды. Сколько килограммов изюма получится из $60$~кг винограда?\par
  \vspace{1.5mm}\hfill\textbf{\color{wcprimary}Ответ:}~\AnswerField{5}%
}

\WorksheetQR{T3-mix/L1/v1}

\end{document}
